% ============================================================================
% 第 3 章 - QEMU 裸机 AArch64
% ============================================================================
\chapter{QEMU 裸机 AArch64}

用户态汇编跑在 macOS 内核之上，依赖系统调用完成 I/O。
本章我们把这个环境彻底去掉——直接跑在一个"没有任何操作系统"的
虚拟 ARM 机器上。

这件事看似疯狂，其实非常简单：你写的汇编程序不再依赖 \texttt{sys\_write}、
\texttt{sys\_read}，而是直接\textbf{读写特定内存地址}。
那个地址上挂了一个 UART 串口设备，你写什么，它就通过串口把字节送到终端；
它收到什么，就读回给你。

整个流程：

\begin{enumerate}
    \item 用 clang（--target=aarch64-none-elf）把 \texttt{start.s} 汇编成 .o
    \item 用 ld.lld 把 .o 链接到地址 \texttt{0x4000\_0000}（QEMU virt 机器的默认载入点）
    \item 启动 \texttt{qemu-system-aarch64}，让它把我们的 ELF 文件载入到
        该地址，然后 CPU 从 \_start 开始执行
\end{enumerate}

\section{start.s：最简裸机程序}
% ============================================================================
% 3.1 start.s
% ============================================================================
\subsection{极简的裸机程序}

这是一个完整的"在裸机上打印一句话然后永远等待"的程序。

\begin{lstlisting}[style=asm,caption=QEMU/start.s — 最简裸机程序]
.section .text
.global _start

_start:
    ldr x0, =0x09000000      ; UART0 的基地址（QEMU virt 机器的约定）
    ldr x1, =msg              ; 字符串所在的地址

print_loop:
    ldrb w2, [x1], #1        ; 读 x1 指向的字节，然后 x1++
    cbz  w2, halt             ; 读到 0（字符串结尾）就退出

wait_uart:
    ldr  w3, [x0, #0x18]     ; 读 UART 的 FLAG 寄存器 (UARTFR, 偏移 0x18)
    tbnz w3, #5, wait_uart   ; bit 5 = TXFF (发送 FIFO 满)，忙等直到不忙
    str  w2, [x0]            ; 把字节写到 UART 的数据寄存器 (UARTDR, 偏移 0x0)
    b    print_loop

halt:
    wfe                      ; 等待事件：进入低功耗状态
    b    halt                ; 万一被唤醒，继续等待

.section .rodata
msg:
    .asciz "Hello QEMU AArch64 bare metal!\n"
\end{lstlisting}

这不到 20 行做了下面几件事：

\paragraph{步骤 1：把两个关键地址取到寄存器}
\begin{itemize}
    \item \texttt{x0 = 0x0900\_0000} —— QEMU 模拟的 ARM PL011 UART0 的
        MMIO 基地址。只要往这个地址写数据，QEMU 就会把它当成字符显示到终端。
    \item \texttt{x1 = msg} —— 字符串在程序内存中的起始地址。
\end{itemize}

\paragraph{步骤 2：逐字节打印}
每一个字符的处理都要先\textbf{查询 UART 是否空闲}：

\begin{itemize}
    \item 读 \texttt{[x0, \#0x18]}（UARTFR 寄存器），检查第 5 位 (TXFF)。
    \item TXFF = Transmit FIFO Full。为 1 表示发送缓冲区已满，
        CPU 必须等待。
    \item 用 \texttt{tbnz w3, \#5, wait\_uart} 忙等这一位变 0。
    \item 一旦空闲，\texttt{str w2, [x0]} 把当前字节写出去。
\end{itemize}

PL011 是 ARM 官方定义的 UART 硬件规范，QEMU 按这个标准实现了
一个虚拟版本。所以你只要按照 PL011 的寄存器表操作，就能驱动它。

\paragraph{步骤 3：停机}
程序打印完字符串（遇到 \texttt{0} 字节即 \texttt{.asciz} 的字符串尾）后，
执行 \texttt{wfe}（Wait For Event）。这是一条 ARM 指令，
让 CPU 停止取指、进入低功耗状态，直到外部中断唤醒它。
我们没有任何中断源被启用，所以它实际上就是永久停机了。

\subsection{和 ISA 章节的差异}

与 ISA/ 目录里的用户态汇编相比，这段代码有几个不同之处：

\begin{itemize}
    \item 没有 ADRP/ADD 组合——因为我们不要求位置无关代码。
        链接器会把 \_start 放到 \texttt{0x4000\_0000}，
        所以直接用 \texttt{ldr reg, =address} 就能拿到绝对地址。
    \item 没有 \texttt{svc \#0x80}——我们直接操作硬件，没有内核帮你做事。
    \item 入口是 \texttt{\_start} 不是 \texttt{\_main}——链接脚本
        指定入口点为 \texttt{\_start}，直接在这里开始执行。
    \item \texttt{.section .rodata} 把只读数据明确放在只读段里。
\end{itemize}

整个程序唯一真正"执行"的部分就是打印——它不会从用户键盘读，
也不会做更复杂的事。后面你想扩展它的话，就是在 \texttt{halt} 之前
加更多功能（比如轮询 UART 的接收寄存器 \texttt{[x0, \#0x00]} 从用户键盘读）。


\section{run.sh：交叉编译与启动脚本}
% ============================================================================
% 3.2 run.sh：交叉编译与启动
% ============================================================================
\subsection{为什么需要三步式构建}

在 macOS 上写一个在 QEMU 内跑的 ARM64 裸机程序，
你面对的是"跨平台"问题：
你的电脑是 ARM64 架构，但你写的程序跑在另一个"ARM64 机器"里。
实际上，它们\textbf{指令集是一样的}，区别只在地址空间和系统调用约定上。

具体来说：

\begin{itemize}
    \item macOS 的 AArch64 用户态程序通过 \texttt{svc \#0x80} 调用内核
    \item 裸机 AArch64 程序根本没有内核，也没有 svc 约定
    \item 因此我们需要一个"不链接任何系统库"的纯裸机二进制
\end{itemize}

因此构建是三步：

\begin{lstlisting}[style=bashstyle,caption=QEMU/run.sh — 编译脚本]
#!/bin/zsh
set -e
mkdir -p bin

# 第一步：用 clang 作为交叉汇编器
# --target=aarch64-none-elf 告诉 clang 目标是裸机 ARM64
#  -c 只编译不链接，产出 bin/start.o
/opt/homebrew/opt/llvm/bin/clang \
    --target=aarch64-none-elf \
    -c start.s \
    -o bin/start.o

# 第二步：用 LLVM 的 ld.lld 链接
# -Ttext=0x40000000 告诉链接器把 .text 段放到 0x4000\_0000
#    这是 QEMU virt 机器的内核默认加载地址
# -e _start 指定入口符号
/opt/homebrew/opt/lld/bin/ld.lld \
    -Ttext=0x40000000 \
    -e _start \
    -o bin/kernel.elf bin/start.o

# 第三步：启动 QEMU
# -machine virt ：使用 virt 虚拟机器（通用 ARM64 虚拟平台）
# -cpu cortex-a72 ：模拟 ARM Cortex-A72 处理器
# -nographic ：用当前终端作为串口控制台（不弹出图形窗口）
# -kernel bin/kernel.elf ：把 ELF 文件作为内核加载
qemu-system-aarch64 \
    -machine virt \
    -cpu cortex-a72 \
    -nographic \
    -kernel bin/kernel.elf
\end{lstlisting}

\paragraph{地址 0x4000\_0000 的来源}
QEMU 的 virt 机器把"内核可加载"的 RAM 起始地址固定为 \texttt{0x4000\_0000}
（1 GB 处）。你编译时让链接器把 \texttt{\_start} 放到这个地址，
然后 QEMU 从这个地址开始执行你的程序。

如果你想让程序执行在别的地址，可以随意改这个数字，
但必须保证它在 QEMU 的 DRAM 范围内（0x4000\_0000 到 0x8000\_0000 之间是一个安全的默认区）。

\paragraph{退出 QEMU}
在 \texttt{-nographic} 模式下，\textbf{Ctrl-A X}（先按 Ctrl-A，松开，再按 X）
可以退出 QEMU。Ctrl-C 是被送进模拟机器里的，不会让你退出 QEMU 主程序。

\subsection{常见错误}

\begin{itemize}
    \item \textbf{command not found: ld.lld}：你的 Homebrew 没装 lld。
        用 \texttt{brew install lld} 或直接 \texttt{brew install llvm} 装完整 LLVM。
    \item \textbf{No 'kernel' option found for this machine type}：QEMU 版本太老
        或你没装 AArch64 版的 QEMU。确认你装的是 \texttt{qemu-system-aarch64}。
    \item \textbf{程序跑了但没输出}：检查 UART 基地址是否为 \texttt{0x09000000}。
        不同机器类型使用不同地址（例如 \texttt{versatilepb} 用的是 \texttt{0x101f1000}）。
    \item \textbf{地址异常或崩溃}：链接地址与 QEMU 加载地址不一致。
        用 \texttt{-Ttext=0x40000000} 配合 \texttt{-kernel bin/kernel.elf}，
        会按链接地址原样加载各段到对应物理地址。
\end{itemize}

运行脚本后，终端会输出：

\begin{verbatim}
Hello QEMU AArch64 bare metal!
\end{verbatim}

然后程序进入 \texttt{wfe} 永久等待。按 Ctrl-A X 退出 QEMU。


\section{理解地址空间与 UART}
% ============================================================================
% 3.3 地址空间与 UART
% ============================================================================
\subsection{裸机地址空间的样子}

QEMU 的 virt 机器把物理内存和外设都映射到同一个地址空间里。
CPU 可以通过普通的 \texttt{ldr/str} 指令直接访问它们。
这叫\textbf{内存映射 I/O（MMIO）}。

\begin{center}
\begin{tabular}{l l}
\hline
地址范围 & 内容 \\
\hline
\texttt{0x0000\_0000} -- \texttt{0x0800\_0000} & Flash / 只读区（QEMU 放了固件）\\
\texttt{0x0900\_0000} -- \texttt{0x0900\_0FFF} & \textbf{UART0 串口}（我们使用这个）\\
\texttt{0x0901\_0000} -- \texttt{0x0901\_0FFF} & UART1 串口 \\
\texttt{0x4000\_0000} -- \texttt{0x8000\_0000} & \textbf{DRAM（1 GB）}：程序放在这里 \\
\texttt{0x0800\_0000} -- \texttt{0x0800\_0FFF} & GIC（中断控制器）\\
\hline
\end{tabular}
\end{center}

所以当你执行 \texttt{str w2, [0x0900\_0000]} 时，
CPU 发出的这个写请求落在 DRAM 之外——它被路由到 UART0 设备，
而不是写进内存。硬件（QEMU 模拟的硬件）截获这个地址，把 \texttt{w2}
的值当作一个要发送的字节。

同样地，执行 \texttt{ldr w3, [0x0900\_0018]} 时，
CPU 从 UART 的 FLAG 寄存器读取当前状态位。

\subsection{PL011 UART 的寄存器表（关键部分）}

ARM PL011 是一个标准的串口控制器，QEMU 精确模拟了它。
每一个寄存器都是 32 位（4 字节）宽：

\begin{center}
\begin{tabular}{l l l}
\hline
偏移 & 名称 & 说明 \\
\hline
\texttt{0x00} & UARTDR（数据寄存器） & 写 = 发送一个字节；读 = 读最近收到的字节 \\
\texttt{0x18} & UARTFR（标志寄存器） & bit 4 = RXFE（接收空）；bit 5 = TXFF（发送满）\\
\hline
\end{tabular}
\end{center}

其他寄存器（波特率、线路控制等）在 QEMU 的 PL011 里通常已经被
模拟层初始化为正常可用的状态，所以最小程序不必设置它们。
你只要能向 DR 发送字节、检查 FR 是否忙碌即可。

下面我们逐行解释 \texttt{start.s} 中最关键的两条指令：

\paragraph{读标志寄存器 — \texttt{ldr w3, [x0, \#0x18]}}
\texttt{x0} 指向 \texttt{0x09000000}，加偏移 \texttt{0x18} 后
得到 FLAG 寄存器地址。第 5 位是"发送 FIFO 满"标志，若为 1 就必须等待。

\paragraph{\texttt{tbnz w3, \#5, wait\_uart} — 忙等直到不忙}
Test Bit and Branch if Not Zero：如果 \texttt{w3} 的第 5 位是 1，
就跳回 \texttt{wait\_uart} 标签继续读。否则（发送器空闲），
继续执行。这是硬件层面的"忙等"循环——真实硬件里这个过程可能需要
几十纳秒；QEMU 里几乎瞬间完成。

\paragraph{\texttt{str w2, [x0]} — 发送字节}
把 \texttt{w2} 的内容写入 \texttt{0x09000000}（UARTDR）。
QEMU 看到这个地址被写入，就把低 8 位当作一个 ASCII 字符送到终端。

\subsection{如果要从用户键盘接收数据}

发送的反方向操作是：
\begin{enumerate}
    \item 读 \texttt{UARTFR}（0x09000018），检查第 4 位（RXFE = Receiver FIFO Empty）
    \item 如果 RXFE=0，说明接收到了新字节
    \item 读 \texttt{UARTDR}（0x09000000），取出低 8 位，就是用户敲的字符
\end{enumerate}

这对应第 4 章 Verilog 里的 UART 接收器做的事——它们本质是同一套
硬件逻辑，一个由你用 Verilog 设计，一个由 QEMU 用 C 模拟。

\subsection{内存与硬件设备的区别}

理解这个模型最重要的一点是：

\begin{center}
\textbf{地址本身没有意义。把地址映射到 RAM 还是设备，由硬件（或 QEMU）决定。}
\end{center}

\begin{itemize}
    \item 写 \texttt{0x4000\_0000} — 写到 DRAM，下次读会读到你写的值。
    \item 写 \texttt{0x0900\_0000} — 写到 UART，下次读这个地址也许读到的是一个刚接收到的字节。
    \item 读 \texttt{0x0900\_0018} — 读到的不是 RAM，是硬件状态寄存器。
\end{itemize}

在真实芯片里，这是由 SOC（System on Chip）内部的\textbf{地址解码器}
实现的——它根据地址范围把读写请求路由到不同的模块。
你在第 4 章里做的 Verilog UART，就是这个模型中的一个模块。


\section{硬件资源概览}
% 硬件资源

目前使用和计划使用的硬件资源：

\begin{itemize}
    \item \textbf{CPU}：Cortex-A72（AArch64），启动时位于 EL2
    \item \textbf{UART}：PL011，MMIO 基址 \texttt{0x09000000}
    \item \textbf{GIC}：Generic Interrupt Controller，用于中断路由
    \item \textbf{Timer}：ARM Generic Timer，用于调度器时钟
    \item \textbf{RAM}：从 \texttt{0x40000000} 开始
    \item \textbf{virtio-net}：后续用于网络功能
\end{itemize}


\section{内存管理}
% ============================================================================
% 内存管理
% ============================================================================

\subsection{内存管理}

我们在裸机环境中实现了一个简单的内存分配器，支持分配和释放，还支持块合并。

内存布局：
- 每个内存块有 8 字节头部：
  - 4 字节：块大小
  - 4 字节：使用标记（0 未使用，1 已使用）
- 堆大小：4096 字节

内存分配器有三个函数：
\begin{itemize}
    \item \texttt{memory\_init}：初始化堆
    \item \texttt{memory\_alloc}：分配内存
    \item \texttt{memory\_free}：释放内存
\end{itemize}

\subsubsection{汇编实现}

\lstinputlisting[style=asm]{../QEMU/src/ASM/memory.s}

\subsubsection{C 语言封装}

C 语言中有对汇编函数的声明：

\lstinputlisting[style=cstyle]{../QEMU/src/include/memory.h}

\subsubsection{工作原理}
\begin{enumerate}
    \item \texttt{memory\_init}：初始化堆，第一个块大小设为 4088 字节（4096 - 8 字节头部）
    \item \texttt{memory\_alloc}：首次适配算法，找到足够大的空闲块
    \item \texttt{memory\_free}：释放内存并尝试与相邻空闲块合并
\end{enumerate}


\section{异常处理}
% ============================================================================
% 异常处理
% ============================================================================

\subsection{异常处理}

在 AArch64 架构中，异常是一个重要的机制。我们实现了完整的异常向量表和异常处理程序。

\subsubsection{异常向量表}

异常向量表要求 2048 字节对齐，包含 16 个入口：

\lstinputlisting[style=asm]{../QEMU/src/ASM/exceptions.s}

\subsubsection{异常处理 C 函数}

我们用 C 语言实现了异常解析和处理：

\lstinputlisting[style=cstyle]{../QEMU/src/exceptionhandler.c}

\subsubsection{\texttt{ESR\_EL2} 寄存器}

\texttt{ESR\_EL2} 寄存器存储异常原因，\texttt{EC} 字段（bit[31:26]）表示异常类别：

\begin{itemize}
    \item \texttt{0x0E}：非法执行状态
    \item \texttt{0x15}：SVC 指令（AArch64）
    \item \texttt{0x21}：指令中止（当前 EL）
    \item \texttt{0x25}：数据中止（当前 EL）
    \item 等等...
\end{itemize}

\subsubsection{工作流程}
\begin{enumerate}
    \item 异常发生，CPU 跳转到异常向量表
    \item 汇编保存所有通用寄存器
    \item 读取 \texttt{ESR\_EL2}、\texttt{ELR\_EL2}、\texttt{FAR\_EL2} 寄存器
    \item 调用 C 函数 \texttt{c\_exception\_handler} 解析和打印
    \item 如果是同步异常，返回 1 表示需要跳过触发指令
    \item 汇编恢复寄存器，执行 \texttt{eret} 返回
\end{enumerate}


\section{后续路线图}
% 后续路线图

以下是从当前基础往完整微型内核推进的路径：

\begin{enumerate}
    \item 物理内存分配器 Page Allocator
    \item 协作式内核线程，手动 yield 切换上下文
    \item 定时器中断与抢占式调度
    \item MMU 与虚拟内存
    \item 系统调用与进程，通过 SVC 指令实现用户态与内核态切换
    \item virtio-net 网络栈
\end{enumerate}

